%==============================================================================
% 13_configuration_clutter.tex — Clutter/blocker foundations for the SSP/JGP
% grouping structure (Thread D, with Prof. Wagler).
%
% STATUS (2026-07-02, Claude-Fable draft): foundations only. The propositions are
% elementary and proved inline; the questions are open. AI-assisted working
% scratch, NOT ground truth — verify before citing or building on externally.
% Machine checks: the star witness and non-matroid witness are hand-verifiable in
% one line each; both also appear in _verification/VERIFIED_FACTS.md (conflict-
% graph corrections block).
%==============================================================================
\documentclass[11pt,a4paper]{article}
\usepackage{amsmath,amssymb,amsthm}
\usepackage[margin=1.1in]{geometry}
\usepackage[numbers]{natbib}
\newtheorem{proposition}{Proposition}
\newtheorem{definition}{Definition}
\newtheorem{question}{Question}
\newtheorem{remark}{Remark}
\newcommand{\Kstar}{K^{\ast}}
\newcommand{\ZSSP}{Z^{\ast}}

\title{The grouping structure of the SSP as an independence system:\\
clutter, blocker, and the right polyhedral questions}
\author{Working note (Thread D)}
\date{\today}

\begin{document}
\maketitle

\section{The group complex and its circuits}

Fix an SSP/JGP instance: jobs $J$, tool sets $T_j$, capacity $b$, $U=\bigcup_j
T_j$.

\begin{definition}[Group complex]
Let $\mathcal{F}=\{S\subseteq J:\ \lvert\bigcup_{j\in S}T_j\rvert\le b\}$, the
family of feasible groups. $\mathcal{F}$ is nonempty (singletons, as
$\lvert T_j\rvert\le b$) and downward closed, hence an \emph{independence
system} (abstract simplicial complex) on ground set $J$. Its \emph{circuits}
$\mathcal{C}$ are the inclusion-minimal infeasible sets; by minimality
$\mathcal{C}$ is a \emph{clutter}. The JGP value $\Kstar$ is the minimum number
of faces of $\mathcal{F}$ covering $J$.
\end{definition}

\begin{proposition}[$\Kstar$ is a hypergraph chromatic number]\label{prop:chrom}
A set of jobs is feasible iff it contains no circuit; hence $\Kstar$ equals the
chromatic number of the circuit hypergraph $(J,\mathcal{C})$ (minimum number of
colours with no monochromatic circuit). The conflict graph $G_{\mathrm{conf}}$
is exactly the family of \emph{two-element} circuits, so
$\chi(G_{\mathrm{conf}})\le\Kstar$, and the inequality is strict in general:
for $T_1=\{0,1\}$, $T_2=\{0,2\}$, $T_3=\{0,3\}$, $b=3$, all pairwise unions have
size $3$ (no conflicts, $\chi=1$) while the triple union has size $4$, so
$\{1,2,3\}$ is a $3$-element circuit and $\Kstar=2$.
\end{proposition}
\begin{proof}
A feasible set contains no infeasible subset (downward closure), in particular
no circuit; conversely a set containing no circuit is feasible, since any
infeasible set contains a minimal infeasible subset. Colour classes of a proper
colouring of $(J,\mathcal{C})$ are thus exactly feasible groups, giving the
identification. Two-element circuits are precisely the conflicting pairs;
colourings of $(J,\mathcal{C})$ refine colourings of $G_{\mathrm{conf}}$, giving
the inequality, and the star computes as stated.
\end{proof}

\begin{proposition}[$\mathcal{F}$ is not a matroid]\label{prop:nonmatroid}
The exchange axiom fails in general: for $T_1=\{1\}$, $T_2=\{2\}$,
$T_3=\{3,4\}$, $b=2$, both $\{3\}$ and $\{1,2\}$ index feasible groups, but
neither $\{3,1\}$ nor $\{3,2\}$ is feasible (unions of size $3$). Hence no
matroid machinery (rank axioms, greedy optimality, matroid-union covering
theorems) applies to $\mathcal{F}$ without further hypotheses.
\end{proposition}
\begin{proof}
Immediate: $\lvert T_3\rvert=2\le b$, $\lvert T_1\cup T_2\rvert=2\le b$,
$\lvert T_3\cup T_i\rvert=3>b$ for $i=1,2$.
\end{proof}

\begin{remark}[What is already known to fail]
Perfectness of $G_{\mathrm{conf}}$ does not predict the heuristic gap: the
$6$-ring's conflict graph is the triangular prism (perfect) with gap $1$, the
$5$-ring's is $C_5$ (imperfect) with gap $0$ (see 05, verified). Any polyhedral
criterion for the gap must therefore see more than the two-element circuits ---
by Proposition~\ref{prop:chrom}, at least the full circuit clutter.
\end{remark}

\section{The covering side: matrices, LP integrality, blocker}

Let $M$ be the job--group incidence matrix (rows $J$, columns the
inclusion-\emph{maximal} feasible groups). JGP is the set-covering program
$\min\{\mathbf 1^\top x: Mx\ge\mathbf 1,\ x\in\{0,1\}^{\cdot}\}$, whose linear
relaxation is empirically very strong (column-generation solves it with small
gaps \citep{Crama1994Column,Otiai2026}).

\begin{question}[Idealness / integrality]\label{q:ideal}
For which instance classes is the covering polyhedron
$\{x\ge0: Mx\ge\mathbf 1\}$ integral --- equivalently, when does the clutter of
minimal covers behave ideally in the sense of \citet{Cornuejols2001}? A
characterisation would explain \emph{when} the JGP relaxation is tight, rather
than only observing that it usually is. Natural first classes: the $b=3$ edge
family (jobs $=$ graph edges; groups $=$ edge sets spanning $\le3$ vertices),
and Helly-type hypotheses on the tool sets under which
$\chi(G_{\mathrm{conf}})=\Kstar$.
\end{question}

\begin{question}[Blocker and lower bounds]\label{q:blocker}
The blocker of the clutter of minimal covers consists of the minimal
\emph{obstructions}: minimal job sets meeting every cover, i.e.\ witnesses
forcing $\Kstar\ge k$. Do the ring/prism instances (the extremal gap witnesses
of 05) carry a distinguished blocker structure --- and is membership of a
$2$-covered blocker element what forces a bridge tool, connecting the blocker
directly to the re-insertion count $R$ and to the exact $\Kstar=3$ formula
$H=q+\min(x_{ab},x_{bc},x_{ca})$ of 05?
\end{question}

\begin{proposition}[The circuit clutter does not determine the gap ---
ANSWERED 2026-07-02]\label{prop:notclutter}
%% Decided by directed enumeration: _verification/verify_clutter_q3_and_genK.py
%% (66 isomorphism classes over edge-m=5 / mixed b=3 / mixed b=4 families;
%% exactly one class is mixed, witnesses below re-verified exactly).
The heuristic gap is \emph{not} a function of the circuit clutter
$\mathcal{C}$, even together with $(b,\lvert U\rvert,n)$. Witnesses ($b=4$,
$U=\{0,\dots,6\}$, four jobs, $q=3$):
\[
  I_0:\ T=\bigl(\{0,2,6\},\{1,2,3\},\{2,3,4,5\},\{3,4,5\}\bigr),
  \qquad
  I_1:\ T=\bigl(\{0,1,3,6\},\{1,4,5,6\},\{2,4\},\{3,5\}\bigr).
\]
Both instances have the same circuit clutter up to relabelling (every job pair
conflicts except the last two; no circuits of size $\ge3$), and moreover the
same $\Kstar=3$ and the same free-initial $\ZSSP=3$; yet $I_0$ has gap $0$ and
$I_1$ has gap $1$.
\end{proposition}

\begin{remark}[Consequence]
No purely clutter-level (a fortiori no conflict-graph-level) criterion can
decide zero-gap: the invariant must read \emph{tool-level} structure. This is
consistent with, and explains the necessity of, the exact $\Kstar=3$ expression
$H=q+\min\min(x_{ab},x_{bc},x_{ca})$ of 05, whose overlap sizes $x_{uv}$ are
precisely the tool-level data the clutter forgets. The polyhedral programme
with Prof.~Wagler should therefore target the \emph{covering} side
(Question~\ref{q:ideal}, where the clutter is the right object for JGP
tightness), not a clutter-only gap criterion.
\end{remark}

\section{The blocking pair and its LP duality (preparation for Prof.~Wagler)}
%% Added 2026-07-15 after Shankar relayed that Prof. Wagler asked for the
%% blocker/clutter analysis "for the dual", posed on the ring instance. The
%% reconstruction below is the standard blocking-duality reading of that
%% request; the ring computations are verified (hand dual + CBC, k=5..9).

Fix an instance and let $\mathcal{G}$ be the (finite) set of inclusion-maximal
feasible groups. Two clutters pair up by blocking duality:

\begin{definition}[Cluster clutter and its blocker]
On ground set $\mathcal{G}$, let
$\mathcal{A}=\{A_j: j\in J\}$ (minimalised), where
$A_j=\{G\in\mathcal{G}: j\in G\}$ is the \emph{cluster} of job $j$. A
transversal of $\mathcal{A}$ is a set of groups meeting every $A_j$ --- i.e.\
a \emph{cover} of $J$ --- so the blocker $b(\mathcal{A})$ is the clutter of
minimal group-covers, and $\Kstar=\min\{|B|:B\in b(\mathcal{A})\}$.
\end{definition}

\noindent The associated LP pair is the JGP covering relaxation and its dual
packing:
\[
  \tau^*=\min\Bigl\{\textstyle\sum_G x_G:\ \sum_{G\ni j}x_G\ge1\ \forall j,\
  x\ge0\Bigr\}
  \;=\;
  \nu^*=\max\Bigl\{\textstyle\sum_j y_j:\ \sum_{j\in G}y_j\le1\ \forall G,\
  y\ge0\Bigr\},
\]
and $\mathcal{A}$ is \emph{ideal} exactly when $\tau^*=\Kstar$ for all
non-negative weightings; by Lehman's theorem idealness holds for
$\mathcal{A}$ iff it holds for $b(\mathcal{A})$ --- the covering and cluster
sides stand or fall together. This is the precise sense in which the blocker
is ``the dual''.

\begin{proposition}[Circuit--group blocking identity]\label{prop:circblocker}
On ground set $J$, the blocker of the circuit clutter $\mathcal{C}$ is the
clutter of minimal complements of maximal feasible groups:
$b(\mathcal{C})=\{J\setminus G:\ G\in\mathcal{G}\}^{\min}$.
\end{proposition}
\begin{proof}
$S$ meets every circuit iff $J\setminus S$ contains no circuit iff
$J\setminus S$ is feasible; $S$ is a minimal such set iff $J\setminus S$ is
maximal feasible.
\end{proof}

\begin{proposition}[Odd rings are the first non-ideal family]\label{prop:oddring}
On the $k$-ring with $b=3$ (jobs $=$ edges of $C_k$; maximal groups $=$
adjacent edge pairs), $\tau^*=k/2$ while $\Kstar=\lceil k/2\rceil$: the
covering LP is tight for even $k$ and has gap $\tfrac12$ for every odd $k$.
\end{proposition}
\begin{proof}
Primal: weight $\tfrac12$ on each of the $k$ adjacent-pair groups covers each
edge exactly once, so $\tau^*\le k/2$. Dual: $y_j=\tfrac12$ is feasible (each
group contains at most two jobs), so $\nu^*\ge k/2$; LP duality gives
$\tau^*=k/2$. The integer optimum is $\lceil k/2\rceil$.
(Verified computationally for $k=5,\dots,9$.)
\end{proof}

\begin{remark}[Why this matters, and what to ask]
(i) The JGP set-covering relaxation is the engine of the whole
branch-and-price line and is reported empirically tight
\citep{Crama1994Column,Otiai2026}; \cref{prop:oddring} exhibits its
integrality gap already on the $5$-ring --- \emph{odd circular structure is
the first obstruction family}, in exact parallel to odd holes in perfection
theory. (ii) Our gap-side result (\cref{prop:notclutter}: the clutter cannot
decide the SSP heuristic gap) and this covering-side result split the
programme cleanly: \textbf{the clutter is the wrong invariant for the gap but
the right one for JGP tightness}. Questions for the session: is
$\mathcal{A}$ ideal precisely when no odd-ring-like structure appears as a
minor? Does Lehman's machinery (or the MFMC property) characterise the
instance classes where Otiai-style branch-and-price closes at the root? And
does the width--length inequality on $(\mathcal{A},b(\mathcal{A}))$ say
anything about $\Kstar-1$ versus $q$ --- the two lower bounds of the gap
theory?
\end{remark}

%% 2026-07-16: prop:circblocker + prop:oddring (blocking pair, odd-ring
%% non-idealness) are now ALSO in the report, Section ssec:clutter. This note
%% remains the working home for Q1/Q2.
\paragraph{Next steps.} (1) \emph{Done:} the gap-from-clutter question is
settled negatively (Proposition~\ref{prop:notclutter}). (2) \emph{Done in
first form:} the blocking pair is set up with the odd-ring non-idealness
(\cref{prop:oddring}) as the opening result for the Wagler session; extend
the computation across the enumerated families to map where $\tau^*<\Kstar$.
(3) Bring Question~\ref{q:blocker}, the witnesses above, and
\cref{prop:oddring} to the session; the prism/$C_5$ pair, the $I_0/I_1$ pair,
and the odd rings are the instructive examples.

\bibliographystyle{plainnat}
\bibliography{references}
\end{document}
